% !TeX root = ../MAgent_Theory_2.0.tex
\chapter{Curvature Costs and Supplied Hierarchies}
\label{ch:curvature-hierarchy}

Integrating out contents can couple dictionaries even when the model contains no loop prior. The resulting action is the negative log evidence of the observed records at fixed dictionaries. In a Gaussian tangent model its quadratic part charges coded reinterpretations; in the weak-anchor or small-noise limit it charges exactly those that changes of contents cannot absorb. Its linear part depends on the records, so this calculation does not imply that flat dictionaries are always preferred. A second construction restricts a supplied hierarchy to rigid clusters. It gives a parent density under explicit normalization hypotheses and a Gaussian minimizer limit, without selecting the clusters or establishing inherited responses and dynamics.

\section{Contents, records, and supplied parents}
\label{it3:ch-setting}

Work in one linear sector with reference content space $\R^K$. The graph has finite vertex set $\mathsf V$ and directed edge set $\mathsf E$. An edge $e:j\to i$ has receiver-owned dictionary generator $\omega_e\in\mathfrak{gl}(K,\R)$, code $\bar T_e:\R^K\to H_e$, and positive-definite noise covariance $R_e$ on the finite-dimensional message space $H_e$. Fix operating contents $c_i$ and put $\kappa_i\omega=\omega c_i$. The contents $y_i\in\R^K$ have independent self laws $\mathcal N(m_i,P_i^{-1})$ with $P_i$ positive definite. These self laws summarize the declared local models and records in this Gaussian example. They are not general definitions of an agent.

The comparison records and global posterior energy are
\begin{equation}
o_e=\bar T_e(y_i-y_j-\kappa_i\omega_e)+\varepsilon_e,
\qquad\varepsilon_e\sim\mathcal N(0,R_e),
\end{equation}
\begin{equation}
\begin{aligned}
E(y,\omega)&=\frac12\sum_i\Vert y_i-m_i\Vert_{P_i}^2
+\frac1{2\sigma^2}\sum_e\Vert\omega_e\Vert_F^2\\
&\quad+\frac12\sum_e\Vert o_e-\bar T_e(y_i-y_j-\kappa_i\omega_e)\Vert_{R_e^{-1}}^2,
\end{aligned}
\label{it3:energy}
\end{equation}
where $\sigma>0$ and the noises are independent. The matrix anchor is isotropic for the declared reference inner product. Coordinate changes carry this metric with the other constitutive data, so the quantitative statements below hold in the reference presentation in which the anchor is isotropic, and are invariant under the reference changes that preserve it, as in \cref{it1:tangent-model}. There are no grounded dictionary factors and no explicit loop factor: the loop weight is $\lambda=0$. The tangent approximation drops the terms of second and higher order in $(\omega_e,y_i-c_i)$ in $\bar T_e(e^{-\omega_e}y_i-y_j)$, among them the product $\omega_e(y_i-c_i)$. Its operating contents remain fixed. The approximation is consistent only while the posterior spread of $y_i-c_i$ is small compared with $|c_i|$; \cref{it3:weak-anchor-determinant} records what happens when it is not.

Let $Y=(\R^K)^{\mathsf V}$, $X=\mathfrak{gl}(K,\R)^{\mathsf E}$, and $H=\bigoplus_e H_e$. Define
\begin{equation}
(\delta y)_e=\bar T_e(y_i-y_j),\qquad
(\mathcal V\omega)_e=\bar T_e\omega_ec_i,\qquad
P=\bigoplus_iP_i,\qquad R=\bigoplus_eR_e.
\end{equation}
The record term is $\frac12\Vert o-\delta y+\mathcal V\omega\Vert_{R^{-1}}^2$. The map $\delta:Y\to H$ is the coboundary of the code sheaf: the two restrictions to the edge stalk $H_e$ are the same code $\bar T_e$. Its first cohomology on the graph, which has no two-cells, is $H^1=H/\operatorname{im}\delta$ \citep[Section~2.2]{HansenGhrist2019}. Cellular sheaves of this kind are used for discourse by \citet{HansenGhrist2021}; that paper works with $H^0$.

For a directed closed walk $\gamma$, let $(B\omega)_\gamma=\sum_{e\in\gamma}\omega_e$, counting traversals with multiplicity. The corresponding product of dictionaries $\bar\theta_e=e^{\omega_e}$ is ordered along the walk, and $B\omega$ is its first-order term at the identity. Opposite attending edges remain independent; traversing a particular transport in geometric reverse would instead use its inverse. Write $\Delta_e=y_i-y_j$. Then $\sum_{e\in\gamma}\Delta_e=0$ on a directed closed walk.

\definitionheading{A supplied hierarchy}{it3:hierarchy}
Partition the child agents into clusters $I$, introduce a parent $A$ for each cluster, and give it a self law $\mathcal N(m_A,P_A^{-1})$. Add the cross-scale comparison factor
\begin{equation}
\exp\left(-\frac{\tau}{2}\Vert y_i-y_A-\omega_{iA}c_A\Vert^2\right),
\qquad i\in I,\qquad \tau>0,
\label{it3:top-down}
\end{equation}
with an isotropic dictionary anchor of variance $\sigma^2$. This is the tangent form of the conditional law $y_i\mid y_A,\omega_{iA}\sim\mathcal N(e^{\omega_{iA}}y_A,\tau^{-1}I)$ at $y_A=c_A$. It is a comparison factor with identity code and fixed operating vector $c_A$. Edges between different child clusters are called bridges. The partition, parents, and cross-scale factors are supplied model data.

\section{The action induced by contents}
\label{it3:ch-induced}

\remarkheading{Evidence gives an induced action for general laws}{it3:general-induced}
Let $\pi(d\omega)$ be a dictionary prior and $p(dy)$ a content law on standard Borel spaces, with contents and dictionaries independent a priori, so that the joint prior is $p\otimes\pi$. Let $\lambda_o$ be a sigma-finite measure on the record space that dominates every record law, and let $\ell(o\mid y,\omega)$ be a jointly measurable version of the record density. Suppose, for the registered record $o$,
\begin{equation}
Z(\omega)=\int\ell(o\mid y,\omega)p(dy),\qquad
0<\int Z(\omega)\pi(d\omega)<\infty.
\end{equation}
Bayes' theorem for this dominated model gives, for $\lambda_o$-almost every record, the posterior
\begin{equation}
\pi(d\omega\mid o)
=\frac{Z(\omega)\pi(d\omega)}{\int Z(\omega')\pi(d\omega')}.
\end{equation}
Where $0<Z(\omega)<\infty$, the contents induce the action $-\log Z(\omega)$. At one registered record the formula defines a version of the posterior: two jointly measurable versions of the density that differ on a null set of records can give different posteriors at that record. The induced action is therefore defined for a fixed version, canonically the version continuous in $o$ when it exists, as in \cref{it3:induced}. This identity requires no Gaussian family. If the record kernels use the dictionaries only through coded discrepancies, that dependence is inherited by $Z$. Neither monotonicity in a loop defect nor invariance under arbitrary re-readings of contents follows: the content law and the registered records also enter the integral.

\theoremheading{The Gaussian action and its cohomological quadratic part}{it3:induced}
In the model of \cref{it3:energy}, integrating out $y$ leaves the complete dictionary action
\begin{equation}
S_{\mathrm{marg}}(\omega)
=\frac1{2\sigma^2}\Vert\omega\Vert^2
+\frac12\Vert o+\mathcal V\omega-\delta m\Vert_{C_P^{-1}}^2
+\mathrm{const},\qquad
C_P=R+\delta P^{-1}\delta^T.
\label{it3:affine-action}
\end{equation}
Thus the dictionary marginal is Gaussian with precision
\begin{equation}
\Lambda_\omega=\sigma^{-2}I+\mathcal V^TC_P^{-1}\mathcal V.
\label{it3:marginal}
\end{equation}
Within this tangent model, if $P=\epsilon\bar P$ with fixed positive-definite $\bar P$ and $\epsilon\downarrow0$, its record precision tends to $\mathcal V^T\Pi_H\mathcal V$, where
\begin{equation}
\Pi_H=R^{-1}-R^{-1}\delta(\delta^TR^{-1}\delta)^+\delta^TR^{-1}.
\label{it3:projection}
\end{equation}
The superscript $+$ denotes the Moore--Penrose inverse. If instead $R=\rho\bar R$ with fixed positive-definite $\bar R$, fixed $P$ and $\rho\downarrow0$, then $\rho$ times the record precision tends to $\mathcal V^T\Pi_H(\bar R)\mathcal V$, with $\Pi_H(\bar R)$ formed from $\bar R$ in place of $R$. For every $\omega$,
\begin{equation}
(\mathcal V\omega)^T\Pi_H\mathcal V\omega
=\min_{a\in Y}\Vert\mathcal V\omega-\delta a\Vert_{R^{-1}}^2.
\label{it3:quotient}
\end{equation}
This is the squared quotient norm of the class of $\mathcal V\omega$ in $H^1$, the squared norm of its harmonic representative in the $R^{-1}$-weighted discrete Hodge decomposition in degree one \citep[Theorem~3.1]{HansenGhrist2019}, whose graph form goes back to \citet{Eckmann1944}; see also \citet{Lim2020}. It vanishes exactly when
$\bar T_e\omega_ec_i=\bar T_e(a_i-a_j)$ for one $a\in Y$ on every edge. If every receiver has $c_i\neq0$, the rank of the limiting quadratic form on $X$ is
\begin{equation}
\sum_e\operatorname{rank}\bar T_e-K|\mathsf V|+\dim H^0,
\label{it3:rank-general}
\end{equation}
where $H^0=\ker\delta$; for codes of full row rank this is $\dim H^1$.

If every code is injective and the underlying undirected multigraph is connected, the rank is $K(|\mathsf E|-|\mathsf V|+1)$; within the full-row-rank class of \cref{it2:model} injective codes are square and invertible. Suppose in addition that all operating contents equal one $c\ne0$ and the chains of the declared directed closed walks span the real cycle space of that multigraph. Opposite attending edges count as separate edges of it. Then the limiting quadratic part vanishes exactly when
$(B\omega)_\gamma c=0$ for every declared walk. For positive $P$, an absorbable reinterpretation obeys
\begin{equation}
(\delta a)^TC_P^{-1}\delta a\leq a^TPa.
\label{it3:anchor-bound}
\end{equation}
\begin{proof}
Under the content prior, $\delta y+\varepsilon$ has mean $\delta m$ and covariance $C_P$. The registered value in that Gaussian experiment is $o+\mathcal V\omega$. Its density gives \cref{it3:affine-action}, including the affine term, and hence \cref{it3:marginal}. The determinant of $C_P$ is independent of $\omega$ in this tangent model.

Set $A=R^{-1/2}\delta\bar P^{-1/2}$. Then
$C_P^{-1}=R^{-1/2}(I+\epsilon^{-1}AA^T)^{-1}R^{-1/2}$.
On the range of $A$ the middle factor tends to zero; on its orthogonal complement it is identity. The limit is \cref{it3:projection}, the residual quadratic form for weighted least squares. For the small-noise limit write $\bar A=\bar R^{-1/2}\delta P^{-1/2}$; then $\rho C_P^{-1}=\bar R^{-1/2}\rho(\rho I+\bar A\bar A^T)^{-1}\bar R^{-1/2}$, and $\rho(\rho I+\bar A\bar A^T)^{-1}$ tends to the orthogonal projection onto the complement of the range of $\bar A$, which is the same as the range of $\bar R^{-1/2}\delta$. Minimizing over $a$ proves \cref{it3:quotient} and its zero criterion. The kernel of $\Pi_H$ is $\operatorname{im}\delta$. When every $c_i\neq0$, $\omega_e\mapsto\omega_ec_i$ is onto $\R^K$, so $\operatorname{im}\mathcal V=\bigoplus_e\operatorname{im}\bar T_e\supseteq\operatorname{im}\delta$, and the rank of $\mathcal V^T\Pi_H\mathcal V$ is $\dim\operatorname{im}\mathcal V-\dim\operatorname{im}\delta$, which is \cref{it3:rank-general}.

For injective codes and a common $c$ that criterion becomes $\omega_ec=a_i-a_j$. A vector-valued edge cochain has this form exactly when it pairs to zero with every cycle. To see the converse directly, choose a spanning tree and define $a_i$ by signed sums along its paths from a root. Each remaining edge closes a fundamental cycle, whose zero sum gives its required difference. The cycle-spanning hypothesis makes the declared directed-walk tests sufficient for these signed cycle tests; it does not identify an opposite attending edge with geometric reversal. For injective codes on a connected graph $\dim H^0=K$ and $\operatorname{rank}\bar T_e=K$, so \cref{it3:rank-general} is $K|\mathsf E|-K(|\mathsf V|-1)$.

Finally, Gaussian elimination, or completion of the square, gives
$u^TC_P^{-1}u=\min_y[y^TPy+\Vert u-\delta y\Vert_{R^{-1}}^2]$.
Set $u=\delta a$ and evaluate the right side at $y=a$ to obtain \cref{it3:anchor-bound}.
\end{proof}

The cycle-spanning hypothesis is stronger than connectedness. Strong connectivity, with all directed closed walks admitted, is sufficient: a strictly positive circulation is obtained by summing directed cycles containing each edge. Adding a sufficiently large multiple of that circulation to any signed cycle makes it nonnegative. Both nonnegative circulations decompose into directed cycles, so their difference lies in the directed-cycle span. More generally, with all directed closed walks admitted, the directed cycles span the cycle space exactly when the condensation multigraph, obtained by contracting each strong component and keeping the edges between components with multiplicity, is a tree. Directed cycles lie inside strong components and span the cycle space of each \citep{GleissLeydoldStadler2003}, and the cycle rank of the whole graph exceeds the sum of those of its strong components by the cycle rank of the condensation.

For a connected graph where the hypothesis fails, take scalar edges $1\to2$, $2\to3$, $1\to3$, unit codes and noise, $c=1$, and $\mathcal V\omega=(0,0,1)$. There is no nonconstant directed closed walk, but
\begin{equation}
\delta y=(y_2-y_1,y_3-y_2,y_3-y_1),\qquad
\min_a\Vert(0,0,1)-\delta a\Vert^2=\frac13.
\label{it3:dag-example}
\end{equation}
Indeed, $(1,1,-1)$ spans the orthogonal complement of $\operatorname{im}\delta$, and the squared projection onto it is $1/3$. The coded obstruction exists even though this directed graph has no loop product that detects it.

The theorem describes the Hessian of an evidence cost, not a record-independent preference for flatness. With weak anchors the full record action tends to
$\frac12(o+\mathcal V\omega)^T\Pi_H(o+\mathcal V\omega)$.
At zero records it charges the cohomology class of the coded reinterpretation. At nonzero records it can prefer a nonzero class that compensates the observed residual. For finite anchors, the shift by $\delta m$ in \cref{it3:affine-action} must also be retained. At finite anchors the quadratic part also charges absorbable reinterpretations, bounded by \cref{it3:anchor-bound}: one scalar edge with unit code, noise and precisions has $H^1=0$ and record precision $c^2/3$. Only the limits are cohomological. The precision $\Lambda_\omega$ and the coupling of dictionaries through shared contents belong to the exact dictionary marginal. Under the agent-factorized global recognition of \cref{ch:link-dynamics} the optimal recognition covariance has no cross-owner blocks, and only the stationary means carry the induced action.

\remarkheading{The omitted determinant at weak anchors}{it3:weak-anchor-determinant}
The weak-anchor limit of \cref{it3:induced} is a statement inside the declared tangent model, and in that limit the tangent model ceases to describe the record it linearizes. Take injective codes on a connected graph. The common mode $y_i=u$ for all $i$, which spans $\ker\delta$, is constrained by no record, so its posterior precision is $\epsilon\iota^T\bar P\iota$, where $\iota u=(u,\ldots,u)/\sqrt{|\mathsf V|}$. The spread of $y_i-c_i$ is then of order $\epsilon^{-1/2}$, and the dropped product $\omega_e(y_i-c_i)$ is not small. For fixed dictionaries the record that keeps it, $\bar T_e(e^{-\omega_e}y_i-y_j)+\varepsilon_e$, is still linear in $y$, with matrix $L(\omega)$ and $L(0)=\delta$; it is the coboundary of the code sheaf twisted by the dictionaries. Gaussian integration of the contents is therefore exact at fixed $\omega$, and the evidence is the expression of \cref{it3:open-fluctuation} with the dictionary-dependent content precision $J(\omega)=\epsilon\bar P+L(\omega)^TR^{-1}L(\omega)$. Since $L(\omega)\iota u=-\mathcal V_u\omega/\sqrt{|\mathsf V|}+O(|\omega|^2)$, with $(\mathcal V_u\omega)_e=\bar T_e\omega_eu$, a Schur complement on the common mode gives
\begin{equation}
\tfrac12\log\det J(\omega)=\mathrm{const}+\tfrac12\log\det\bigl(\epsilon\iota^T\bar P\iota+\mathcal G(\omega)\bigr)+r(\omega),\qquad
\mathcal G(\omega)_{uu'}=\frac{(\mathcal V_u\omega)^T\Pi_H(\mathcal V_{u'}\omega)}{|\mathsf V|},
\label{it3:determinant-crossover}
\end{equation}
where $\mathcal G$ is read as a quadratic form in $u,u'\in\R^K$, higher-order terms in $\omega$ are omitted, and $r(\omega)$ has curvature bounded as $\epsilon\downarrow0$. For $|\omega|^2\ll\epsilon$ the dictionary dependence is $\frac1{2\epsilon}\Tr\bigl((\iota^T\bar P\iota)^{-1}\mathcal G(\omega)\bigr)$, of order $\epsilon^{-1}$, whereas the tangent quadratic part stays bounded. This term charges $\omega$ whenever $\mathcal V_u\omega\notin\operatorname{im}\delta$ for some content direction $u$, including directions with $\omega c=0$: with $K=2$, $\bar T_e=I_2$, $c=e_1$ and $\omega_e=E_{22}$ on every edge of a graph with a directed cycle, the tangent penalty vanishes while $\mathcal G(\omega)\neq0$. At weak anchors the omitted determinant therefore dominates the tangent quadratic part. In the small-noise limit $R=\rho\bar R$ at fixed anchors both terms scale as $\rho^{-1}$, and the same Schur computation, with $P$ in place of $\epsilon\bar P$, makes the ratio of the determinant term to the tangent term of the order of the common-mode prior variance divided by $|c|^2$. The cohomological part is therefore the leading dictionary dependence of the evidence in the window noise variance $\ll$ content prior variance $\ll|c|^2$, which is also the regime in which the tangent approximation is consistent.

\remarkheading{Which experiment sees a dictionary}{it3:observation}
The matter term at one fixed operating content is blind to a perturbation $A$ when $\bar T_eAc_i=0$. Grounded messages with a referent law of positive-definite second moment instead test $\bar T_eA$, as their means range over the referents. These are different kernels. With $K=2$, $\bar T=I_2$, $c=(1,0)^T$, and $A=E_{22}$, the first experiment sees no change, whereas grounded referents identify $A$. A loop prior can constrain directions left free by a particular matter experiment, but those directions need not be untestable in another admitted experiment.

There is a further distinction between law comparisons and link observables. For a loop map $H$ and law $q$, $\KL(q\Vert H_\#q)$ vanishes whenever $H$ preserves $q$. For positive-definite $\Sigma$, the covariance-dressed link defect
\begin{equation}
Q(H;\Sigma)=\frac12\Vert\Sigma^{-1/2}(I-H)\Sigma^{1/2}\Vert_F^2
\label{it3:link-defect}
\end{equation}
vanishes exactly at $H=I$, since the similarity transform is invertible. A quarter-turn preserves $q=\mathcal N(0,I_2)$, giving zero law divergence, while $Q(H;I_2)=2$. The choice of defect therefore does not inherit the law stabilizer automatically.

A shared-sample relay around a loop is one way to probe its action. Another is to identify the constituent dictionaries from sufficiently informative grounded edge records and compose them. The evidence calculation requires neither direct loop observations nor such an identification procedure. Its conclusion is that records couple dictionaries through shared contents in the declared joint.

\section{Stationary reinterpretation}
\label{it3:ch-stationary}

\propositionheading{One operating content gives a rank-one stationary mean}{it3:rank-one}
At a stationary point of \cref{it3:energy}, each dictionary satisfies
\begin{equation}
\omega_e=\sigma^2\bar T_e^T
\bigl(R_e+\sigma^2|c_i|^2\bar T_e\bar T_e^T\bigr)^{-1}
(\bar T_e\Delta_e-o_e)c_i^T.
\label{it3:rank-one-eq}
\end{equation}
It has rank at most one. For registered agreement, $o_e=0$,
\begin{equation}
\omega_ec_i=F_e\Delta_e,\qquad
F_e=\sigma^2|c_i|^2\bar T_e^T
\bigl(R_e+\sigma^2|c_i|^2\bar T_e\bar T_e^T\bigr)^{-1}\bar T_e.
\label{it3:gain}
\end{equation}
These are also equations for the stationary Gaussian recognition means. In the plug-in per-agent model the same equation holds with $\Delta_e=\mu_i-\mu_j$. In its law-level version, with peer covariance $\Sigma_j$ fixed at the considered state, replace $R_e$ by $\widetilde R_e=R_e+\bar T_e\Sigma_j\bar T_e^T$. For the tangent reciprocal potential at fixed positive covariances and positive row weights $\beta_e$, use $o_e=0$ and replace $R_e$ by $(R_e+\bar T_e\Sigma_j\bar T_e^T)/\beta_e$. In all cases the operating contents are fixed, the anchor is isotropic, and there are no loop or grounded terms. The replacements do not assert equality of dictionary forces across these objectives.
\begin{proof}
Only the anchor and the edge record depend on $\omega_e$. Its normal equation is
\begin{equation}
(\sigma^{-2}I+\kappa_i^T\bar T_e^TR_e^{-1}\bar T_e\kappa_i)\omega_e
=\kappa_i^T\bar T_e^TR_e^{-1}(\bar T_e\Delta_e-o_e).
\end{equation}
The adjoint is $\kappa_i^Tv=vc_i^T$ for the Frobenius inner product, and $\kappa_i\kappa_i^T=|c_i|^2I$. Applying the push-through identity gives \cref{it3:rank-one-eq}; multiplication by $c_i$ gives \cref{it3:gain}. A quadratic Gaussian free energy has the same mean normal equation as its energy. In the per-agent objectives only the owner's record and anchor contain this dictionary, with the stated convolution covariance in the law-level case. In the reciprocal objective its mean-dependent edge cost is $\frac{\beta_e}{2}\Vert\bar T_e(\mu_i-\mu_j-\omega_ec_i)\Vert_{(R_e+\bar T_e\Sigma_j\bar T_e^T)^{-1}}^2$, which gives the last replacement.
\end{proof}

\corollaryheading{Equal gains cancel the linearized loop sum}{it3:dissonance}
Let the operating content be a common $c\ne0$, register $o=0$, and suppose the stationary equations of \cref{it3:rank-one} hold, with the effective noise of the chosen objective. Then
$(B\omega)_\gamma=g_\gamma c^T$ for a vector $g_\gamma$, and
\begin{equation}
(B\omega)_\gamma c
=\sum_{e\in\gamma}F_e\Delta_e
=\sum_{e\in\gamma}(F_e-\bar F)\Delta_e
\label{it3:holonomy-gain}
\end{equation}
for every fixed matrix $\bar F$. Thus the linearized loop map annihilates $c^\perp$. The exact ordered product has the same property. Every stationary dictionary has the form $\omega_e=v_ec^T$; since $(vc^T)^n=(c\cdot v)^{n-1}vc^T$, one has $e^{vc^T}=I+\varphi(c\cdot v)vc^T$ with $\varphi(t)=(e^t-1)/t$, and $(I+uc^T)(I+hc^T)=I+(h+u(1+c\cdot h))c^T$. Hence $H_\gamma=I+h_\gamma c^T$ exactly, where $h$ is updated along the walk by $h\leftarrow h+u_e(1+c\cdot h)$ with $u_e=\varphi(c\cdot v_e)v_e$, and the residual dissonance of the loop is the single vector $h_\gamma=(H_\gamma c-c)/|c|^2$. It vanishes when all gains on the loop agree, or when every disagreement on the loop is zero. On a two-cycle $e:j\to i$, $e':i\to j$, its action on $c$ is $(F_e-F_{e'})\Delta_e$.
\begin{proof}
Every term of \cref{it3:rank-one-eq} has right factor $c^T$. Summing the equations and applying them to $c$ gives the first equality in \cref{it3:holonomy-gain}. The second uses $\sum_\gamma\Delta_e=0$. A matrix $gc^T$ with $c\ne0$ is zero exactly when its action on $c$ is zero, which proves the cancellation statements.
\end{proof}

Nonzero stationary loop action therefore requires both disagreement and variation of the gains in this fixed-topic tangent experiment; those conditions are not sufficient. The gains depend on code, noise, and anchor strength. If all self means coincide and all records register agreement, the strictly convex global energy has the common-content, zero-dictionary minimizer. Conflicting self means or other registered records change this conclusion. For general laws the evidence depends on the complete content distributions, and one operating vector need not determine a stationary dictionary. Even within the Gaussian example, cancellation of $B\omega$ is a first-order statement: equal gains make the generator sum vanish but need not make the second-order part of $h_\gamma$ vanish.

\section{Rigid clusters and their boundary dictionaries}
\label{it3:ch-hierarchy}

\remarkheading{Reinterpretation limits the effective top-down precision}{it3:cap}
Integrating out the free Gaussian cross-scale dictionary in \cref{it3:top-down} leaves a factor in $y_i-y_A$ with precision
\begin{equation}
(\tau^{-1}+\sigma^2|c_A|^2)^{-1}I.
\end{equation}
This follows because $\omega_{iA}c_A$ has covariance $\sigma^2|c_A|^2I$ and adds to the cross-scale record noise. For $c_A\ne0$ the precision is bounded above by $(\sigma^2|c_A|^2)^{-1}I$ as $\tau$ increases. For $c_A=0$ it equals $\tau I$ and there is no such finite cap. Pinning the cross-scale dictionaries removes this source of reinterpretation. This elimination identity alone does not prove that disagreement concentrates on a cluster boundary.

\propositionheading{A normalized rigid restriction of a supplied hierarchy}{it3:rigid}
Let child and parent contents lie in one standard Borel space $E$ with a common sigma-finite reference measure $\mu$, and let $G$ be a measurable dictionary group with identity $I$. Specify measurable self densities $p_i,p_A$ with respect to $\mu$. Let the record likelihoods have the form $\ell_e(o_e\mid d_e(z_i,z_j,\theta_e))$, where $d_e:E\times E\times G\to H_e$ is a declared measurable discrepancy with $d_e(z,z,I)=0$ and $\ell_e(o_e\mid\cdot)$ is Borel on $H_e$. Fix internal and cross-scale dictionaries at the identity. Suppose every internal or cross-scale likelihood at its registered record and zero discrepancy is a finite positive constant. For every supplied cluster require
\begin{equation}
0<Z_A=\int_E p_A(z)\prod_{i\in I}p_i(z)\mu(dz)<\infty.
\label{it3:rigid-normalizer}
\end{equation}
Evaluate the joint density on $y_i=y_A$ and use $\mu$ once for each parent coordinate, together with a declared bridge-dictionary prior $\pi_b$. If the dictionary prior couples bridge dictionaries with pinned internal or cross-scale ones, as the anchored loop prior of \cref{it2:model} does around loops containing both, $\pi_b$ is the normalized restriction of the joint prior density at the pinned values, assumed to have a finite positive normalizer; for the Gaussian loop prior this restriction retains the wall terms $\frac\lambda2\Vert\sum_{e\in\gamma\cap\mathsf E_b}\omega_e\Vert_{W_\gamma}^2$, where $\mathsf E_b$ is the set of bridges, which couple bridge dictionaries. Assume that the bridge factors, integrated against the parent self laws \cref{it3:parent-density} and $\pi_b$, have a finite positive integral, $0<\int\prod_{\mathrm{bridges}}\ell_e\prod_Ap_A'(z_A)\mu(dz_A)\pi_b(d\theta)<\infty$. Then the resulting normalized rigid restriction is a communication model with parent self density
\begin{equation}
p_A'(z)=Z_A^{-1}p_A(z)\prod_{i\in I}p_i(z)
\label{it3:parent-density}
\end{equation}
and one record factor for each original bridge, retaining its code, dictionary, and registered value. The parent density is a product of the child and parent self densities, a logarithmic pool with unit weights \citep{GenestZidek1986}.

For the Gaussian tangent family, with $\mu$ Lebesgue measure on $\R^K$, its parent self parameters are
\begin{equation}
P_A'=P_A+\sum_{i\in I}P_i,\qquad
P_A'm_A'=P_Am_A+\sum_{i\in I}P_im_i.
\label{it3:parent-gaussian}
\end{equation}
A bridge retains its fixed operating vector $c_i$; these vectors can be bridge dependent at the parent level unless they are common within each receiving cluster. With positive-definite self precisions, identity internal and cross-scale dictionaries, and the Gaussian cross-scale factors of precision $\tau$, the normalized soft Gaussian law converges to the rigid restriction with Lebesgue measure, with mean and covariance errors $O(\tau^{-1})$; in particular the global stationary minimizer tends to the minimizer of this rigid parent model at that rate. A reference measure $g(z)dz$ instead multiplies the parent density by $g^{-|I|}$, so \cref{it3:parent-gaussian} requires constant $g$, and the finiteness of $Z_A$ can depend on that choice.
\begin{proof}
On equal contents an internal discrepancy is zero, and the hypotheses replace each internal and cross-scale record by a finite positive constant. The self factors at a cluster multiply at their common argument. Dividing by $Z_A$ gives \cref{it3:parent-density}; the assumed remaining normalizer makes the entire parent law a probability law. Each bridge retains its original likelihood evaluated at the two parent contents. In the tangent model this argument is $\bar T_e(y_{A(i)}-y_{A(j)}-\omega_ec_i)$. Completing the square in the Gaussian self factors gives \cref{it3:parent-gaussian}.

For the minimizer limit collect all free Gaussian variables, including bridge dictionaries, in $x$. The energy is
$E_\tau(x)=\frac{\tau}{2}\Vert Cx\Vert^2+\frac12x^THx-h^Tx+\mathrm{const}$,
where $Cx$ lists the child-parent differences and $H$ is positive definite. Choose orthonormal bases $U$ for $\ker C$ and $V$ for its orthogonal complement, and write $x=Ua+Vb$. With $H_{aa}=U^THU$ and the analogous blocks, the first normal equation gives
$a=a_0-H_{aa}^{-1}H_{ab}b$, where $a_0=H_{aa}^{-1}U^Th$ is the constrained minimizer. The second gives
\begin{equation}
\left[\tau V^TC^TCV+H_{bb}-H_{ba}H_{aa}^{-1}H_{ab}\right]b
=V^Th-H_{ba}a_0.
\end{equation}
Both the first coefficient and the Schur complement are positive definite on the $b$ space. Hence $b=O(\tau^{-1})$ and $a-a_0=O(\tau^{-1})$. The same block decomposition of the precision $\tau C^TC+H$ gives its inverse as $UH_{aa}^{-1}U^T+O(\tau^{-1})$, and $Ua_0$ and $UH_{aa}^{-1}U^T$ are the mean and covariance of the Gaussian law restricted to $\ker C$ with Lebesgue measure there, which is the rigid restriction because $\ker C$ is parameterized linearly by the parent contents and the bridge dictionaries. This proves the claims.
\end{proof}

The reference measure and normalizers are part of this restriction. It is neither conditioning a continuous law on a probability-zero diagonal without a specified convention nor integrating out children. Individually normalized self densities do not guarantee \cref{it3:rigid-normalizer}: two copies of $p(z)=\frac13z^{-2/3}$ on $(0,1)$ have a product proportional to $z^{-4/3}$, whose integral is infinite. Disjoint supports can instead make the product integral zero. The Gaussian convergence statement does not extend by itself to arbitrary laws. In the Gaussian case the normalized cross-scale kernel tends to a Dirac mass with respect to Lebesgue measure, which is why that measure appears; for general laws the cross-scale kernels would have to select the reference measure.

\propositionheading{Flat wall walks and a common boundary dictionary}{it3:flat-wall}
Consider two clusters $I,J$ whose internal dictionaries are identity. Assume that bridges exist in both directions. For every bridge $e:j\to i$ from $J$ to $I$ and every bridge $g:i'\to j'$ from $I$ to $J$, require directed internal paths $i\to i'$ in $I$ and $j'\to j$ in $J$. Internal strong connectivity is sufficient. Every edge of these connecting paths, including any cross-scale edge they use, carries the identity dictionary. Let the bridge dictionaries be arbitrary elements of a group $G$.

Every directed walk that uses $e$, such internal paths, and $g$ has holonomy $\bar\theta_g\bar\theta_e$. All these walks are flat exactly when there is one $\bar\theta_D\in G$ such that all bridges $J\to I$ carry $\bar\theta_D$ and all bridges $I\to J$ carry $\bar\theta_D^{-1}$. At first order the corresponding condition is $\omega_e=D$, $\omega_g=-D$.

For rigid clusters with common operating content $c$ and zero bridge records, if the free bridge dictionaries satisfy \cref{it3:rank-one}, the linearized wall loop acts on $c$ by $(F_e-F_g)(y_I-y_J)$. More generally, with identity internal dictionaries but nonzero internal disagreements, that action is
\begin{equation}
(\omega_e+\omega_g)c
=(F_e-F_g)\Delta_e+F_g(\Delta_e+\Delta_g).
\label{it3:wall-residual}
\end{equation}
\begin{proof}
The internal factors are identity, so the ordered product is $\bar\theta_g\bar\theta_e$. The path hypothesis realizes every opposite bridge pair as a walk of this kind. Fix one bridge in each direction, which is possible by nonemptiness. The equations $\bar\theta_g\bar\theta_e=I$ for all pairs imply that every forward bridge is the inverse of the fixed backward bridge, and every backward bridge is the inverse of that common forward value. The converse follows by multiplication. Replacing products by their first-order sums gives the generator condition. The stationary bridge equations give
$\omega_ec+\omega_gc=F_e\Delta_e+F_g\Delta_g$, which is \cref{it3:wall-residual}; rigidity makes $\Delta_g=-\Delta_e$.
\end{proof}

The identity requirement on the connecting paths is needed: with $I=\{a_1,a_2\}$, $J=\{b\}$, bridges $e:b\to a_1$ with $\bar\theta_e=3$ and $g:a_2\to b$, and a path $a_1\to A\to a_2$ through a parent with scalar dictionaries $2$ and $1/4$, the choice $\bar\theta_g=2/3$ makes the only wall walk flat while $\bar\theta_g\bar\theta_e=2$, so no common $\bar\theta_D$ exists. The all-pairs path hypothesis can be weakened. Form the bipartite pairing graph whose vertices are the bridges and whose edges join $e$ and $g$ when an admissible walk of the stated kind closes them. If this graph is connected and both directions are nonempty, the conclusion holds, by propagating $\bar\theta_g=\bar\theta_e^{-1}$ along its edges; if it is disconnected and $G$ is nontrivial, different components can carry different values of $\bar\theta_D$.

Without the path hypothesis, separate reciprocal bridge pairs can each be flat while carrying different dictionaries. For $I=\{a_1,a_2\}$ and $J=\{b_1,b_2\}$ with no internal directed paths, let $b_1\leftrightarrow a_1$ carry scalar dictionaries $1,1$ and $b_2\leftrightarrow a_2$ carry $2,1/2$. Every available directed loop is flat, but the bridges from $J$ to $I$ disagree. Parent-to-child edges alone supply no return paths between these pairs.

When the proposition applies, all records in one bridge direction read a common dictionary through their own codes. Grounded referents with positive-definite second moment identify it up to the intersection of those code stabilizers; an injective stacked code identifies it completely for the location-family experiments. The group-level flatness statement is independent of Gaussianity. Its stationary gain formula is a Gaussian tangent statement. Neither licenses replacing distinct dictionaries by one merely because a single-topic experiment fails to distinguish them.

\remarkheading{Repeating a supplied rigid restriction}{it3:tower}
The density factorization can be repeated along a supplied hierarchy whenever the required measures and normalizers exist. In the Gaussian family the parent self term pools child information through \cref{it3:parent-gaussian}, while the bridge factors retain the boundary records. At finite coupling, integrating out children instead gives a Schur complement; representing it within the prescribed comparison-record family is an additional closure requirement.

These statements do not select a partition, prove its persistence, or construct an inherited response and update kernel for a predeclared probe class. They also do not establish compression on a post-handoff invariant set, interaction necessity against a null flow, or a second selection depending on the first. Those are separate requirements of the formation standard. A supplied rigid parent density therefore does not by itself establish recursive formation or complete dynamical inheritance.

\section{Remaining extensions}
\label{it3:ch-next}

\openproblemheading{Formation with top-down feedback}{it3:open-formation}
Determine whether a specified dynamics can select persistent blocks and create parents whose feedback supports further selection. The construction must retain the laws, records, evaluators, and complete response and update data required by its admitted protocols, and must meet the compression, null-flow, and counterfactual-recursion conditions. The supplied partitions and minimizer limits above do not settle these questions. This is the specialization of \cref{inh:formation-open} to constructions with declared top-down factors. One candidate in which the link sector does selection work is a wall across a frustrated cut: loops inside two candidate clusters admit flat assignments while cross-cluster loops do not, and the source rows would have to concentrate inside the flat clusters above a loop-weight threshold while the null flow and the zero-loop-weight control form no wall. The record dependence of \cref{it3:induced} decides the sign of such an effect, and inheritance of a learned link sector is open.

\openproblemheading{Fluctuations beyond the fixed tangent experiment}{it3:open-fluctuation}
In the present tangent model $C_P$ is dictionary independent, so its determinant contributes no dictionary force. For a nonlinear link model with Gaussian conditional content precision $J(\omega)$, information vector $b(\omega)$, and scalar energy term $k(\omega)$, Gaussian integration gives
$k(\omega)-\frac12b(\omega)^TJ(\omega)^{-1}b(\omega)+\frac12\log\det J(\omega)$
up to a constant, provided $J(\omega)$ is positive definite. At weak anchors the log-determinant is the leading dictionary dependence by \cref{it3:weak-anchor-determinant}, and it charges the directions with $\omega c=0$. For unit codes and $U(1)$-valued links the content precision is $P$ plus the connection Laplacian of the multigraph of directed links, and Forman's theorem expands its determinant over cycle-rooted spanning forests with nonnegative cycle weights $2-2\cos$ of the holonomies \citep{Forman1993}; \citet{Kenyon2011} extends the expansion to vector bundles. For codes of rank below $K$ the operator is a sheaf Laplacian, for which no such expansion is derived here. Determine the complete action in a declared nonlinear communication model, including its off-topic directions and the signs of the resulting loop terms. The quadratic penalty of \cref{it3:induced} alone does not supply that calculation.

\openproblemheading{General content laws}{it3:open-general}
Identify hypotheses under which the evidence action for general laws admits a useful decomposition into absorbable reinterpretations and residual incompatibility. Such a result must preserve record dependence and distinguish an invariance of the likelihood from an invariance of the content measure. Establishing a normalized soft-to-rigid limit also requires more than the finite product normalizers used for the static restriction; in the Gaussian family it holds with Lebesgue measure by \cref{it3:rigid}, and for general laws the reference measure would be selected by the cross-scale kernels. The stationary rank-one and gain formulas are not claimed outside the Gaussian tangent family.

\openproblemheading{Learning gains from loop records}{it3:open-gains}
For a declared grounded or relay experiment, determine which combinations of code, noise, and dictionary-anchor strength can be inferred from observed loop effects. The gain-weighted formula is an identity at a stationary tangent state; an identifiability theorem requires the full admitted record experiment and cannot be inferred from that identity alone.
